\documentclass[11pt]{article}

\usepackage{amsmath,amssymb,amsthm,mathtools}
\usepackage[margin=1in]{geometry}
\usepackage{microtype}

\newtheorem{theorem}{Theorem}[section]
\newtheorem{lemma}[theorem]{Lemma}
\newtheorem{corollary}[theorem]{Corollary}
\newtheorem{remark}[theorem]{Remark}

\newcommand{\dd}{\,\mathrm d}

\title{Triangle Books with a Page-Rooted Two-Edge Tail:\\
Fractional Tail Convexity}
\author{}
\date{}

\begin{document}

\maketitle

\begin{abstract}
Let $P_m$ be a book of $m\geq1$ triangles on a common spine edge, with
one otherwise disjoint path of length two attached to a page vertex.  We
prove the chromatic-polynomial graphon bound for every $P_m$ and every
$p\geq1/2$.  Symmetrization over the page orbit distributes the rooted
two-edge-path density $A=T_Wd$ with exponent $1/m$.  Conditional
factorization over the pages and Jensen's inequality reduce the problem to
the convex scalar function
\[
 a^{1/m}\max\{2a-p,0\}.
\]
This yields the sharp target directly from $\int A=\int d^2\geq p^2$.
The case $m=2$ proves the formerly open NetworkX Atlas graph $123$.
\end{abstract}


%%%%%%%%%%%%%%%%%%%%%%%%%%%%%%%%%%%%%%%%%%%%%%%%%%%%%%%%%%%%%%%%%%%%%%%%%%%%%%%
\section{The family and its target}
%%%%%%%%%%%%%%%%%%%%%%%%%%%%%%%%%%%%%%%%%%%%%%%%%%%%%%%%%%%%%%%%%%%%%%%%%%%%%%%

Let $(\Omega,\mu)$ be an arbitrary probability space and let
$W\colon\Omega^2\to[0,1]$ be symmetric and measurable.  Put
\[
 p=\int_{\Omega^2}W(x,y)\dd\mu(x)\dd\mu(y),
 \qquad
 d(x)=\int_\Omega W(x,y)\dd\mu(y),
\]
and define
\[
 A(x)=(T_Wd)(x)=\int_\Omega W(x,y)d(y)\dd\mu(y).
\]
The rooted triangle density is
\[
 \tau(x)=\int_{\Omega^2}
 W(x,y)W(x,z)W(y,z)\dd\mu(y)\dd\mu(z).
\]
All integrands used below are bounded, nonnegative, and measurable.

For an integer $m\geq1$, let $P_m$ have spine vertices $a,b$, page
vertices $z_1,\ldots,z_m$, and tail vertices $u,v$.  Its edge set is
\[
 ab,\qquad az_i,bz_i\quad(1\leq i\leq m),
 \qquad z_1u,uv.
\]
Thus the two-edge tail is rooted at one page, not at a spine endpoint.

After the ordered spine is properly coloured, each page has $x-2$ choices.
The two successive tail vertices each have $x-1$ choices.  Consequently
\begin{equation}
 \chi_{P_m}(x)=x(x-1)^3(x-2)^m,
 \qquad \chi(P_m)=3.
 \label{eq:chromatic}
\end{equation}
For $q=1-p$, direct substitution gives the polynomial target
\begin{equation}
 \Phi_{P_m}(p)
 =q^{m+4}\chi_{P_m}(1/q)
 =p^3(2p-1)^m.
 \label{eq:target}
\end{equation}
The required density interval is therefore $p\in[1/2,1]$.


%%%%%%%%%%%%%%%%%%%%%%%%%%%%%%%%%%%%%%%%%%%%%%%%%%%%%%%%%%%%%%%%%%%%%%%%%%%%%%%
\section{Fractional tail convexity}
%%%%%%%%%%%%%%%%%%%%%%%%%%%%%%%%%%%%%%%%%%%%%%%%%%%%%%%%%%%%%%%%%%%%%%%%%%%%%%%

The elementary inequality $rs\geq r+s-1$ for $r,s\in[0,1]$, applied to
$W(x,y),W(x,z)$ and multiplied by $W(y,z)$, gives
\begin{equation}
 \tau(x)\geq2A(x)-p.
 \label{eq:pointwise}
\end{equation}
Since $\tau\geq0$, we may use the stronger combined statement
\begin{equation}
 \tau(x)\geq\max\{2A(x)-p,0\}.
 \label{eq:positive-part}
\end{equation}

\begin{lemma}[Fractional tail convexity]
\label{lem:convexity}
Let $p>0$ and $0<\alpha\leq1$.  The function
\[
 f_{p,\alpha}(a)
 =a^\alpha\max\{2a-p,0\},
 \qquad 0\leq a\leq1,
\]
is convex and nondecreasing.
\end{lemma}

\begin{proof}
The function vanishes on $[0,p/2]$.  On $[p/2,1]$ it equals
\[
 2a^{\alpha+1}-pa^\alpha.
\]
Its first two derivatives on that interval are
\begin{align*}
 f_{p,\alpha}'(a)
 &=a^{\alpha-1}\bigl(2(\alpha+1)a-p\alpha\bigr),\\
 f_{p,\alpha}''(a)
 &=\alpha a^{\alpha-2}
 \bigl(2(\alpha+1)a-p(\alpha-1)\bigr).
\end{align*}
For $a\geq p/2$, the first bracket is at least $p>0$.  The second bracket
is positive because $\alpha\leq1$.  At $a=p/2$, the right derivative is
$p(p/2)^{\alpha-1}>0$, whereas the left derivative is zero.  Thus the
one-sided derivatives increase across the junction, proving both claims.
\end{proof}

\begin{lemma}[Fractionally weighted rooted triangle]
\label{lem:weighted-triangle}
If $p\in[1/2,1]$ and $0<\alpha\leq1$, then
\begin{equation}
 \int_\Omega A(x)^\alpha\tau(x)\dd\mu(x)
 \geq p^{1+2\alpha}(2p-1).
 \label{eq:weighted-triangle}
\end{equation}
\end{lemma}

\begin{proof}
Symmetry and Fubini give
\[
 \int_\Omega A(x)\dd\mu(x)
 =\int_\Omega d(x)^2\dd\mu(x)
 \geq p^2,
\]
where the inequality is Jensen.  By
\eqref{eq:positive-part}, Lemma~\ref{lem:convexity}, Jensen's inequality,
and monotonicity,
\begin{align*}
 \int A^\alpha\tau
 &\geq\int f_{p,\alpha}(A)\\
 &\geq f_{p,\alpha}\!\left(\int A\right)\\
 &\geq f_{p,\alpha}(p^2).
\end{align*}
The condition $p\geq1/2$ is exactly what is needed for $p^2\geq p/2$.
Therefore
\[
 f_{p,\alpha}(p^2)
 =p^{2\alpha}(2p^2-p)
 =p^{1+2\alpha}(2p-1),
\]
which proves \eqref{eq:weighted-triangle}.
\end{proof}


%%%%%%%%%%%%%%%%%%%%%%%%%%%%%%%%%%%%%%%%%%%%%%%%%%%%%%%%%%%%%%%%%%%%%%%%%%%%%%%
\section{The page-tail theorem}
%%%%%%%%%%%%%%%%%%%%%%%%%%%%%%%%%%%%%%%%%%%%%%%%%%%%%%%%%%%%%%%%%%%%%%%%%%%%%%%

\begin{theorem}[Page-rooted two-edge tails on triangle books]
\label{thm:main}
Let $m\geq1$.  Every graphon $W$ of edge density $p\in[1/2,1]$ satisfies
\[
 \boxed{
 t(P_m,W)\geq p^3(2p-1)^m=\Phi_{P_m}(p).
 }
\]
\end{theorem}

\begin{proof}
Integrating the two tail vertices first gives the book integrand multiplied
by $A(x_{z_1})$.  The $m$ pages form one automorphism orbit.  Average the
$m$ representations in which the tail is rooted at page $z_j$ and apply
the arithmetic--geometric mean inequality.  This gives
\begin{align}
 t(P_m,W)
 &\geq
 \int_{\Omega^{m+2}}W(x,y)
 \prod_{j=1}^m
 \left[
 W(x,z_j)W(y,z_j)A(z_j)^{1/m}
 \right]
 \dd\mu^{m+2}.
 \label{eq:symmetrized}
\end{align}
For fixed $x,y$, put
\[
 R(x,y)=\int_\Omega
 W(x,z)W(y,z)A(z)^{1/m}\dd\mu(z).
\]
The page variables in \eqref{eq:symmetrized} are independent, so
\begin{equation}
 t(P_m,W)\geq\int_{\Omega^2}W(x,y)R(x,y)^m\dd\mu^2.
 \label{eq:factorization}
\end{equation}

Here $p\geq1/2>0$.  Normalize $W(x,y)\dd\mu^2/p$ to a probability
measure and apply Jensen to the convex function $r\mapsto r^m$:
\begin{equation}
 \int W R^m
 \geq p^{1-m}\left(\int W R\right)^m.
 \label{eq:edge-jensen}
\end{equation}
Fubini and the definition of $\tau$ give the exact identity
\begin{equation}
 \int_{\Omega^2}W(x,y)R(x,y)\dd\mu^2
 =\int_\Omega A(z)^{1/m}\tau(z)\dd\mu(z).
 \label{eq:rooted-identity}
\end{equation}
Apply Lemma~\ref{lem:weighted-triangle} with $\alpha=1/m$ and substitute
into \eqref{eq:factorization}--\eqref{eq:edge-jensen}:
\begin{align*}
 t(P_m,W)
 &\geq p^{1-m}
 \left[p^{1+2/m}(2p-1)\right]^m\\
 &=p^3(2p-1)^m.
\end{align*}
This is precisely \eqref{eq:target}.
\end{proof}

At $p=1/2$, the target vanishes and every displayed step remains valid;
no division by $2p-1$ occurs.  At $p=1$, $W=1$ almost everywhere and both
sides equal one.  Every balanced complete $k$-partite graphon, $k\geq2$,
realizes equality: $d=p$, $A=p^2$, and $\tau=p(2p-1)$ almost everywhere,
so all symmetrization, factorization, and scalar inequalities are sharp.


%%%%%%%%%%%%%%%%%%%%%%%%%%%%%%%%%%%%%%%%%%%%%%%%%%%%%%%%%%%%%%%%%%%%%%%%%%%%%%%
\section{Atlas 123 and the formalization boundary}
%%%%%%%%%%%%%%%%%%%%%%%%%%%%%%%%%%%%%%%%%%%%%%%%%%%%%%%%%%%%%%%%%%%%%%%%%%%%%%%

\begin{corollary}[Atlas 123]
NetworkX Atlas graph $123$, with graph6 code \texttt{EJy?}, is positive.
It has order six, size seven, chromatic number three, and
\[
 \chi_{F_{123}}(x)=x(x-1)^3(x-2)^2.
\]
Every graphon of edge density $p\in[1/2,1]$ satisfies
\[
 t(F_{123},W)\geq p^3(2p-1)^2=\Phi_{F_{123}}(p).
\]
\end{corollary}

\begin{proof}
Its NetworkX edge list is
\[
 04,\ 05,\ 12,\ 13,\ 14,\ 23,\ 24.
\]
Take spine $12$, pages $3,4$, and tail $4-0-5$.  This is exactly $P_2$,
so Theorem~\ref{thm:main} applies.
\end{proof}

The proof uses ordinary homomorphism densities on an arbitrary probability
space.  All uses of Tonelli--Fubini involve bounded nonnegative integrands.
For formal verification, the reusable analytic statement is
Lemma~\ref{lem:weighted-triangle} for real $0<\alpha\leq1$, built from the
positive-part rooted-triangle bound, fractional scalar convexity, and
$\int A=\int d^2$.  The remaining ingredients are finite orbit averaging,
arithmetic--geometric mean, conditional page factorization, and Jensen under
the normalized edge measure.  The graph-specific module awaiting inspection
is \texttt{lean/Taeyoung/Examples/Graph123.lean}.

The tail must be rooted at a page vertex of a single triangle book.  The
theorem does not cover a tail at a spine endpoint (handled separately in
\texttt{triangle\_book\_two\_edge\_tail.tex}), a two-edge tail on $K_4$
(Atlas 142), or attachments to a general mixed-chordal graph.

The graph, target, fractional exponents, scalar derivatives, and Atlas
identification are reconstructed exactly by
\begin{verbatim}
python codes/verify_triangle_book_page_two_edge_tail.py
\end{verbatim}


%%%%%%%%%%%%%%%%%%%%%%%%%%%%%%%%%%%%%%%%%%%%%%%%%%%%%%%%%%%%%%%%%%%%%%%%%%%%%%%
\section{What the Lean formalization actually proves}
%%%%%%%%%%%%%%%%%%%%%%%%%%%%%%%%%%%%%%%%%%%%%%%%%%%%%%%%%%%%%%%%%%%%%%%%%%%%%%%

The machine-checked development is
\texttt{lean/Taeyoung/Methods/PageTail/\{Core,Rows\}.lean}, ending at
\texttt{satisfiesLowerBound\_123}.  It follows this note except at one place,
and covers only the scoped case $m=2$.

\paragraph{The departure: Lemma~\ref{lem:convexity} is not formalized.}
This note proves that $f_{p,\alpha}(a)=a^\alpha\max\{2a-p,0\}$ is convex and
nondecreasing on $[0,1]$ --- by computing $f'$ and $f''$ on $(p/2,1]$ and
comparing one-sided derivatives at $a=p/2$ --- and then applies Jensen at
$\int A\geq p^2$ to obtain Lemma~\ref{lem:weighted-triangle}.  The
formalization never differentiates anything and never mentions
\texttt{ConvexOn}.  At $\alpha=1/2$ it proves the affine minorant that Jensen
would have produced, directly:
\begin{equation}
 p^2(2p-1)+\Bigl(3p-\tfrac12\Bigr)(a-p^2)
 \leq\sqrt a\,\max\{2a-p,0\}
 \qquad (a\geq0,\ \tfrac12\leq p\leq1),
 \label{eq:lean-tangent-sqrt}
\end{equation}
as \texttt{PageTail.tangent\_sqrt}.  Substituting $a=s^2$ with $s=\sqrt a\geq0$
makes both cases polynomial: for $2a\geq p$ the difference of the two sides is
\[
 2(s-p)^2\Bigl(s+\tfrac p2+\tfrac14\Bigr)\geq0,
\]
a double root at $s=p$ times a positive factor, which is exactly the tangency
this note obtains by differentiating; and for $2a<p$ the right-hand side is
zero while the left-hand side, increasing in $a$ because $3p-\frac12>0$, is at
most its value at $a=p/2$, namely $-p\bigl(p-\frac12\bigr)^2\leq0$.

Integrating \eqref{eq:lean-tangent-sqrt} needs only $\int A\geq p^2$ and the
positive-part bound \eqref{eq:positive-part}, and yields
Lemma~\ref{lem:weighted-triangle} at $\alpha=1/2$, which is
\texttt{PageTail.sqrt\_weighted\_rootedTriangle}.  The general $\alpha$ is not
formalized; the same substitution $a=s^{1/\alpha}$ would not stay polynomial,
so a different certificate would be needed there.

\paragraph{Scope: only $m=2$ is formalized.}
Theorem~\ref{thm:main} is proved here for every $m\geq1$, by orbit averaging,
arithmetic--geometric mean and Jensen for $r\mapsto r^m$.  At $m=2$ each of
those three steps collapses to a Cauchy--Schwarz, and the Lean file does
exactly that:
\[
 R(x,y)^2\leq S(x,y)G(x,y),\qquad
 \Bigl(\int\!\!\int WR\Bigr)^{2}\leq p\int\!\!\int WR^2,
\]
with $R=\int W(x,\cdot)W(y,\cdot)\sqrt A$.  The first is the weighted
Cauchy--Schwarz of \texttt{PureChordal/WeightedCauchySchwarz.lean} applied
inside the page with weight $W(x,\cdot)W(y,\cdot)$; the second is the same
lemma on $\mu\otimes\mu$ with weight $W$.  Nothing is lost for the catalogue:
$P_1$ is the five-vertex triangle with a two-edge tail, already
\texttt{verified}, and $P_3$ has seven vertices.

\paragraph{Everything else transcribes directly.}
The pointwise bound \eqref{eq:pointwise} is \texttt{Link.rootedTriangle\_ge};
\eqref{eq:positive-part} is \texttt{max\_le} against
\texttt{rootedTriangle\_nonneg}; $\int A=\int d^2$ is
\texttt{BookTail.integral\_pathOp}; and \eqref{eq:rooted-identity} is
\texttt{PageTail.integral\_edge\_pageWeightOp}, a weight-parametrized clone of
the existing \texttt{Link.integral\_edge\_pageOp} --- the page operator had to
gain an arbitrary vertex weight, since the page carrying the tail weighs $A$
and the Cauchy--Schwarz that splits it weighs $\sqrt A$.  The final division by
$p$ is \texttt{le\_of\_mul\_le\_mul\_left}, so the argument never divides by
$2p-1$.

\end{document}
